\uuid{EDGQ}
\titre{Théorème de Bézout}
\auteur{bodin ; croizat ; sacre}
\organisation{exo7}
\licence{CC BY-NC-SA 4.0}
\datecreate{2022-03-08}
\niveau{L1}
\module{Algèbre}
\chapitre{Arithmétique dans Z}
\sousChapitre{Nombres premiers, nombres premiers entre eux}
\difficulte{1}
\theme{Théorème de Bézout}
\exo7qcmid{hilisit.arithmetique.017}

\contenu{
\begin{question}
\qclassification{103.04}
 Soit $a$ et $b$ deux entiers tels que $\pgcd(a,b) = 4$. Alors on peut trouver deux entiers $u$ et $v$ tels que :
    \begin{answers}
        \bad{$au-bv=2$}
        \bad{$au+bv=2$}
        \good{$au-bv=4$}
        \good{$au+bv=12$}
    \end{answers}
    \begin{explanations}
    Une égalité $au \pm bv = 2$ nous indiquerait que tout diviseur de $a$ et $b$ diviserait $2$ donc $\pgcd(a,b) = 1$ ou $2$, ce qui n'est pas le cas ici.\\
    Le théorème de Bézout nous garantit l'existence de deux entiers $U$ et $V$ tels que $aU + bV = 4$. Si l'on prend $v=-V$, on obtient $au-bv=4$. Si l'on multiplie l'égalité de Bézout par $3$, on a alors $a \times (3U) + b \times (3V) = 3 \times 4 = 12$.
    \end{explanations}
\end{question}
}
